\hypertarget{chapter-6-macros-as-terms}{%
\chapter{Macros as Terms}\label{chapter-6-macros-as-terms}}

Part I established that TeX processes a document by rewriting tokens rather than executing procedures, and Part II covered how to organize the files a document is built from. This chapter turns to the material an author actually writes when extending LaTeX: macros. Understanding macros correctly, as the rewrite rules described in Chapter 1 rather than as function definitions in a conventional programming language, is what separates writing macros that work by luck from writing macros whose behavior can be predicted before they are ever run.

\hypertarget{def-and-newcommand}{%
\section{\texorpdfstring{\texttt{\textbackslash{}def} and \texttt{\textbackslash{}newcommand}}{\textbackslash def and \textbackslash newcommand}}\label{def-and-newcommand}}

TeX itself provides exactly one primitive for defining a macro: \texttt{\textbackslash{}def}, which takes a control sequence name, an optional parameter text describing how arguments are delimited, and a replacement text. LaTeX's \texttt{\textbackslash{}newcommand} and \texttt{\textbackslash{}renewcommand} are built on top of \texttt{\textbackslash{}def}, adding a small amount of safety: \texttt{\textbackslash{}newcommand} checks that the name being defined is not already in use and raises an error if it is, \texttt{\textbackslash{}renewcommand} checks the opposite, and both provide a more convenient syntax for the common case of a macro taking a fixed number of simple, undelimited arguments, \texttt{\#1} through \texttt{\#9}.

The relationship between the two matters because \texttt{\textbackslash{}def} remains available and is sometimes the better tool. \texttt{\textbackslash{}newcommand} cannot redefine an existing name without \texttt{\textbackslash{}renewcommand}, cannot easily define a macro with a delimited argument (one that reads up to a specific token rather than a fixed number of braces), and adds a small amount of overhead in exchange for its safety checks. A macro author working inside LaTeX's conventions reaches for \texttt{\textbackslash{}newcommand} by default, since the safety it provides catches a real class of accidental redefinition errors, but recognizing when \texttt{\textbackslash{}def}'s lower-level flexibility is actually needed, rather than working around \texttt{\textbackslash{}newcommand}'s limitations with increasingly contorted workarounds, is part of what distinguishes fluent macro writing from macro writing that happens to work.

\hypertarget{expansion-versus-execution}{%
\section{Expansion versus Execution}\label{expansion-versus-execution}}

Chapter 1 introduced the distinction between tokenization and expansion; macro writing is where that distinction has direct, daily consequences. Some TeX operations are expandable: they can occur inside contexts that require full expansion before anything else happens, such as the argument to \texttt{\textbackslash{}edef} or a conditional's test, and their job is to produce more tokens for further processing. Other operations are not expandable at all; they have side effects, such as changing a register's value or starting a new paragraph, and cannot meaningfully occur in a context that only wants tokens.

A macro defined with \texttt{\textbackslash{}def} or \texttt{\textbackslash{}newcommand} is, by default, expandable, in the sense that its name gets replaced by its body wherever expansion reaches it. But the body of that macro is free to contain non-expandable material, in which case the macro as a whole behaves as a mix: expanding partially to reveal its non-expandable content, at which point expansion of that particular token stops, since there is nothing further to expand at that position. This is why some macros can be used inside a context like a section title that gets written to an auxiliary file, and others cannot: the ones that can are built entirely from expandable material all the way down; the ones that cannot contain something, even indirectly, that only makes sense to execute rather than expand.

\texttt{\textbackslash{}edef}, as opposed to \texttt{\textbackslash{}def}, defines a macro whose body is fully expanded once, at definition time, rather than left to be expanded later wherever the macro is used. This is the standard tool for ``freezing'' the current value of something that would otherwise change, such as capturing the current value of a counter into a fixed string rather than a reference that would track the counter's later changes. Confusing \texttt{\textbackslash{}def} and \texttt{\textbackslash{}edef} is one of the most common sources of a macro that behaves correctly the first time it is used and incorrectly the second, because \texttt{\textbackslash{}def} captured a reference to something mutable while \texttt{\textbackslash{}edef} would have captured its value at that moment.

\hypertarget{argument-specification-and-xparse}{%
\section{\texorpdfstring{Argument Specification and \texttt{xparse}}{Argument Specification and xparse}}\label{argument-specification-and-xparse}}

A macro defined with plain \texttt{\textbackslash{}newcommand} takes a fixed number of simple arguments, each either mandatory or, for at most one argument, optional with a default value. Many of the macros an author actually wants to write need more than this: an argument that is optional but has no sensible single default, an argument delimited by a specific character rather than braces, or a variable number of arguments depending on what was provided. The \texttt{xparse} package, part of the LaTeX3 kernel and available for use regardless of whether the rest of a document uses LaTeX3 syntax, provides a compact notation for specifying exactly this kind of argument structure: \texttt{m} for a mandatory argument, \texttt{o} for optional, \texttt{s} for an optional star, \texttt{d} for a delimited argument with custom delimiters, among others, combined into a single specifier string that describes a macro's full calling convention in one place.

This specifier string is worth thinking of as a lightweight signature for the macro: it says, in a compact and checkable form, what a call to this macro is required to look like and what forms are permitted, independent of what the macro's body actually does with the arguments once they are captured. A macro defined with \texttt{\textbackslash{}NewDocumentCommand\{\textbackslash{}keyword\}\{m\ o\}\{...\}} states plainly, at the point of definition, that a call takes one mandatory argument and one optional one, which is more legible than inferring the same fact from reading a \texttt{\textbackslash{}def}-based implementation that manually checks for and strips an optional square-bracket argument by hand. For a macro whose calling convention will be used by anyone other than the person who wrote it, including a future version of the same author, that legibility is worth the small overhead of loading \texttt{xparse}.

\hypertarget{writing-macros-that-predict-their-own-behavior}{%
\section{Writing Macros That Predict Their Own Behavior}\label{writing-macros-that-predict-their-own-behavior}}

The point of treating macros as terms rather than procedures is not merely conceptual accuracy for its own sake; it changes how a macro should be designed. A macro intended for use inside contexts requiring full expansion (section titles, index entries, anywhere text is written to an auxiliary file and reread) has to be built from fully expandable pieces throughout its call chain, not merely at its own top level, since a macro that itself expands cleanly but calls another macro internally that does not will inherit that inner macro's limitation. Checking this by working through what a macro expands to, step by step, in the specific context it will actually be used, catches this class of error before it surfaces as a mysterious failure three chapters and several months later, when the macro is finally used somewhere its original author never tested.

The chapters that follow build directly on this foundation: key-value interfaces, covered next, are a more elaborate argument-handling pattern built from the same primitives introduced here, and robust command design, closing out this part of the book, is largely the discipline of applying the expansion-versus-execution distinction consistently across every macro a project depends on.
